\documentclass[10pt,a4paper]{amsart}
\usepackage{amsmath,amssymb,amsthm,graphicx}
\usepackage{verbatim}
\newenvironment{reference}{\comment}{\endcomment}
%\newenvironment{reference}{}{}
\newenvironment{slogan}{\comment}{\endcomment}
\newenvironment{history}{\comment}{\endcomment}
% For commutative diagrams we use Xy-pic
\usepackage[all]{xy}

% We use 2cell for 2-commutative diagrams.
\xyoption{2cell}
% Theorem environments.
%
\theoremstyle{plain}
\newtheorem{theorem}[subsection]{Theorem}
\newtheorem{proposition}[subsection]{Proposition}
\newtheorem{lemma}[subsection]{Lemma}

\theoremstyle{definition}
\newtheorem{definition}[subsection]{Definition}
\newtheorem{example}[subsection]{Example}
\newtheorem{exercise}[subsection]{Exercise}
\newtheorem{situation}[subsection]{Situation}

\theoremstyle{remark}
\newtheorem{remark}[subsection]{Remark}
\newtheorem{remarks}[subsection]{Remarks}

\numberwithin{equation}{subsection}

% Macros
%
\def\lim{\mathop{\mathrm{lim}}\nolimits}
\def\colim{\mathop{\mathrm{colim}}\nolimits}
\def\Spec{\mathop{\mathrm{Spec}}}
\def\Hom{\mathop{\mathrm{Hom}}\nolimits}
\def\Ext{\mathop{\mathrm{Ext}}\nolimits}
\def\SheafHom{\mathop{\mathcal{H}\!\mathit{om}}\nolimits}
\def\SheafExt{\mathop{\mathcal{E}\!\mathit{xt}}\nolimits}
\def\Sch{\mathit{Sch}}
\def\Mor{\mathop{\mathrm{Mor}}\nolimits}
\def\Ob{\mathop{\mathrm{Ob}}\nolimits}
\def\Sh{\mathop{\mathit{Sh}}\nolimits}
\def\NL{\mathop{N\!L}\nolimits}
\def\CH{\mathop{\mathrm{CH}}\nolimits}
\def\proetale{{pro\text{-}\acute{e}tale}}
\def\etale{{\acute{e}tale}}
\def\QCoh{\mathit{QCoh}}
\def\Ker{\mathop{\mathrm{Ker}}}
\def\Im{\mathop{\mathrm{Im}}}
\def\Coker{\mathop{\mathrm{Coker}}}
\def\Coim{\mathop{\mathrm{Coim}}}

% Boxtimes
%
\DeclareMathSymbol{\boxtimes}{\mathbin}{AMSa}{"02}

%
% Macros for moduli stacks/spaces
%
\def\QCohstack{\mathcal{QC}\!\mathit{oh}}
\def\Cohstack{\mathcal{C}\!\mathit{oh}}
\def\Spacesstack{\mathcal{S}\!\mathit{paces}}
\def\Quotfunctor{\mathrm{Quot}}
\def\Hilbfunctor{\mathrm{Hilb}}
\def\Curvesstack{\mathcal{C}\!\mathit{urves}}
\def\Polarizedstack{\mathcal{P}\!\mathit{olarized}}
\def\Complexesstack{\mathcal{C}\!\mathit{omplexes}}
% \Pic is the operator that assigns to X its picard group, usage \Pic(X)
% \Picardstack_{X/B} denotes the Picard stack of X over B
% \Picardfunctor_{X/B} denotes the Picard functor of X over B
\def\Pic{\mathop{\mathrm{Pic}}\nolimits}
\def\Picardstack{\mathcal{P}\!\mathit{ic}}
\def\Picardfunctor{\mathrm{Pic}}
\def\Deformationcategory{\mathcal{D}\!\mathit{ef}}
\usepackage{cleveref}
\crefname{originalreference}{original reference}{original references}
\Crefname{originalreference}{Original reference}{Original references}
\makeatletter
\newlabel{divisors-lemma-blowing-up-affine}{{divisors-lemma-blowing-up-affine}{source}{}{sourcepointer-1}{}}
\newlabel{divisors-lemma-blowing-up-affine@cref}{{[originalreference][0][]divisors-lemma-blowing-up-affine}{[1][0][]source}{}{}{}}
\makeatother
\crefname{originalreference}{original reference}{original references}
\Crefname{originalreference}{Original reference}{Original references}
\makeatletter
\newlabel{perfect-lemma-quasi-coherence-pullback}{{perfect-lemma-quasi-coherence-pullback}{source}{}{sourcepointer-2}{}}
\newlabel{perfect-lemma-quasi-coherence-pullback@cref}{{[originalreference][0][]perfect-lemma-quasi-coherence-pullback}{[1][0][]source}{}{}{}}
\makeatother
\crefname{originalreference}{original reference}{original references}
\Crefname{originalreference}{Original reference}{Original references}
\makeatletter
\newlabel{coherent-lemma-coherent-on-proj}{{coherent-lemma-coherent-on-proj}{source}{}{sourcepointer-3}{}}
\newlabel{coherent-lemma-coherent-on-proj@cref}{{[originalreference][0][]coherent-lemma-coherent-on-proj}{[1][0][]source}{}{}{}}
\makeatother
\crefname{originalreference}{original reference}{original references}
\Crefname{originalreference}{Original reference}{Original references}
\makeatletter
\newlabel{coherent-lemma-recover-tail-graded-module}{{coherent-lemma-recover-tail-graded-module}{source}{}{sourcepointer-4}{}}
\newlabel{coherent-lemma-recover-tail-graded-module@cref}{{[originalreference][0][]coherent-lemma-recover-tail-graded-module}{[1][0][]source}{}{}{}}
\makeatother
\title{The cohomology sheaves of the derived pullback of a quasi-coherent module on the blowup are annihilated by a fixed power of the ideal}
\author{The Stacks Project Authors}
\date{stacks-project:local-cohomology.tex:local-snapshot-20261012; GNU FDL 1.2 or later}
\begin{document}
\maketitle
\noindent\textit{Editorial note.} All delivered mathematical text is the author's own, copied from the pinned Stacks Project chapter source sources/stacks-local-cohomology/pinned-original-source.tex; no statement or proof step was rewritten, reordered or removed.\par\medskip
\noindent\textit{Editorial note.} Preamble provenance: the Stacks Project keeps one single shared preamble for every chapter; this card takes it directly from the staged sources/stacks-local-cohomology/preamble.tex as three verbatim spans (lines 8-12, 14-18 and 154-217), so no derived file was written and nothing was added or redefined. The wrapper class and the base packages are supplied by the assembler.\par\medskip
\noindent\textit{Editorial note.} Bibliography provenance: neither the counted statement nor its proof contains a citation (this lane rescans the whole span before writing the card), so the card carries no bibliography and no entry was reconstructed.\par\medskip
\noindent\textit{Editorial note.} Selection provenance: this card was opened under the file-based control contract (CONTROL.md), with the mechanical pre-card gate run on its statement and proof interval before the card was written.\par\medskip
\noindent\textit{Allowed standard prerequisites (editorial).}\begin{itemize}
\item The affine open covering of the blowup by the spectra of the affine blowup algebras A[I/a\_i]. Invoked inside the counted proof at source lines 4651-4655 as a cross-chapter pointer. Reduces the statement about sheaves on the blowup to statements about the affine blowup algebras.
\item The description of the derived pullback of the quasi-coherent module on the affine pieces, and the annihilation statement for the affine blowup algebra. The derived pullback description is a cross-chapter pointer (source lines 4656-4658), and the annihilation for the affine case is the author lemma lemma-annihilates-affine, whose statement is reproduced verbatim as a context span of this card (source lines 4554-4560). Turns the required annihilation of the cohomology sheaves into the annihilation of the Tor modules of the affine blowup algebras.
\end{itemize}\medskip
\ifdefined\hypertarget\hypertarget{sourcepointer-1}{}\fi
\noindent\textit{Source reference.} Original reference divisors-lemma-blowing-up-affine: sources/stacks-local-cohomology/pinned-original-source.tex, source lines 4655--4655. Navigation pointer only: the counted text refers to the Stacks item divisors-lemma-blowing-up-affine, which belongs to another chapter of the project and is not part of this card. The pointer carries the target's own label as its number because this chapter's staged snapshot does not contain the other chapter's tag table; its content is not reproduced here.\par\medskip
\ifdefined\hypertarget\hypertarget{sourcepointer-2}{}\fi
\noindent\textit{Source reference.} Original reference perfect-lemma-quasi-coherence-pullback: sources/stacks-local-cohomology/pinned-original-source.tex, source lines 4658--4658. Navigation pointer only: the counted text refers to the Stacks item perfect-lemma-quasi-coherence-pullback, which belongs to another chapter of the project and is not part of this card. The pointer carries the target's own label as its number because this chapter's staged snapshot does not contain the other chapter's tag table; its content is not reproduced here.\par\medskip
\ifdefined\hypertarget\hypertarget{sourcepointer-3}{}\fi
\noindent\textit{Source reference.} Original reference coherent-lemma-coherent-on-proj: sources/stacks-local-cohomology/pinned-original-source.tex, source lines 4612--4612. Navigation pointer only: the counted text refers to the Stacks item coherent-lemma-coherent-on-proj, which belongs to another chapter of the project and is not part of this card. The pointer carries the target's own label as its number because this chapter's staged snapshot does not contain the other chapter's tag table; its content is not reproduced here.\par\medskip
\ifdefined\hypertarget\hypertarget{sourcepointer-4}{}\fi
\noindent\textit{Source reference.} Original reference coherent-lemma-recover-tail-graded-module: sources/stacks-local-cohomology/pinned-original-source.tex, source lines 4613--4613. Navigation pointer only: the counted text refers to the Stacks item coherent-lemma-recover-tail-graded-module, which belongs to another chapter of the project and is not part of this card. The pointer carries the target's own label as its number because this chapter's staged snapshot does not contain the other chapter's tag table; its content is not reproduced here.\par\medskip
\begin{lemma}
\label{lemma-annihilates}
In the situation above, let $t$ be an upper bound on the number of
generators for $I$. There exists an integer $c = c(A, I) \geq 0$
such that for any $A$-module $M$ the cohomology sheaves
$H^j(Lp^*\widetilde{M})$ are annihilated by $I^c$ for $j \leq -t$.
\end{lemma}

\noindent
For the rest of the discussion in this section we fix a Noetherian ring $A$
and an ideal $I \subset A$. We denote
$$
p : X \to \Spec(A)
$$
the blowing up of $\Spec(A)$ in the ideal $I$. In other words, $X$ is the
$\text{Proj}$ of the Rees algebra $\bigoplus_{n \geq 0} I^n$.
By Cohomology of Schemes, Lemmas \ref{coherent-lemma-coherent-on-proj} and
\ref{coherent-lemma-recover-tail-graded-module}
we can choose an integer $q(A, I) \geq 0$ such that for all $q \geq q(A, I)$
we have $H^i(X, \mathcal{O}_X(q)) = 0$ for
$i > 0$ and $H^0(X, \mathcal{O}_X(q)) = I^q$.

\begin{lemma}
\label{lemma-annihilates-affine}
Let $I = (a_1, \ldots, a_t)$ be an ideal of a Noetherian ring $A$.
Set $a = a_1$ and denote $B = A[\frac{I}{a}]$ the affine blowup algebra.
There exists a $c > 0$ such that $\text{Tor}_i^A(B, M)$ is annihilated
by $I^c$ for all $A$-modules $M$ and $i \geq t$.
\end{lemma}

\begin{proof}
Say $I = (a_1, \ldots, a_t)$. The question is affine local on $X$.
For $1 \leq i \leq t$ let $B_i = A[\frac{I}{a_i}]$ be the affine
blowup algebra. Then $X$ has an affine open covering by
the spectra of the rings $B_i$, see
Divisors, Lemma \ref{divisors-lemma-blowing-up-affine}.
By the description of derived pullback given in
Derived Categories of Schemes, Lemma
\ref{perfect-lemma-quasi-coherence-pullback}
we conclude it suffices to prove that for each $i$ there exists a $c \geq 0$
such that
$$
\text{Tor}_j^A(B_i, M)
$$
is annihilated by $I^c$ for $j \geq t$. This is
Lemma \ref{lemma-annihilates-affine}.
\end{proof}

\end{document}
